\begin{lemma}\label{lemma:isParadoxical_bound}\lean{CollatzMapBasics.isParadoxical_bound}\uses{lemma:linear_decomposition'}\leanok
Let $j, n \in \mathbb{N}$ with $n > 0$. If the pair $(j, n)$ is paradoxical, then:
\begin{equation}
    0 < 1 - C(j, n) \le \frac{E(j, n)}{n}.
\end{equation}
\end{lemma}
\begin{proof}\leanok
From the paradoxical condition $C(j, n) < 1$, we get $0 < 1 - C(j, n)$.
By Lemma~\ref{lemma:linear_decomposition'}, $T^j(n) = C(j, n) \cdot n + E(j, n)$.
The paradoxical condition $T^j(n) \ge n$ therefore gives $C(j,n) \cdot n + E(j,n) \ge n$,
hence $E(j,n) \ge (1 - C(j,n)) \cdot n$, and dividing by $n > 0$ yields the result.
\end{proof}

\begin{lemma}\label{lemma:odd_iter_bound_theorem_4_1}\lean{CollatzMapBasics.odd_iter_bound_theorem_4_1}\leanok
Let $j, n, m \in \mathbb{N}$ with $m > 0$. Suppose that every odd iterate $T^k(n)$ with $k < j$ satisfies $T^k(n) \ge m$. Let $q = \mathit{q}(j, n)$ be the number of odd steps in the first $j$ iterations. Then:
\begin{equation}
    T^j(n) \le \frac{n \cdot \left(3 + \frac{1}{m}\right)^q}{2^j}.
\end{equation}
\end{lemma}
\begin{proof}\leanok
By induction on $j$. The base case $j = 0$ is trivial.

For the inductive step, assume the bound holds for $j$. We consider two cases depending on the parity of $T^j(n)$:

\textbf{Odd case} ($X(T^j(n)) = 1$): By definition,
\begin{equation}
    T^{j+1}(n) = T(T^j(n)) \le T^j(n) \cdot \frac{3 + 1/m}{2}.
\end{equation}
Combining with the inductive hypothesis $T^j(n) \le n(3 + 1/m)^q/2^j$ and using $q_{\text{new}} = q + 1$:
\begin{equation}
    T^{j+1}(n) \le \frac{n(3+1/m)^q}{2^j} \cdot \frac{3+1/m}{2} = \frac{n(3+1/m)^{q+1}}{2^{j+1}}.
\end{equation}

\textbf{Even case} ($X(T^j(n)) = 0$): By definition, $T^{j+1}(n) = T^j(n)/2$. Since $q$ stays the same,
\begin{equation}
    T^{j+1}(n) \le \frac{n(3+1/m)^q}{2^j \cdot 2} = \frac{n(3+1/m)^q}{2^{j+1}}. \qedhere
\end{equation}
\end{proof}

\begin{lemma}\label{lemma:correction_ratio_bound_theorem_4_1}\lean{CollatzMapBasics.correction_ratio_bound_theorem_4_1}\uses{lemma:odd_iter_bound_theorem_4_1,lemma:linear_decomposition'}\leanok
Under the same hypotheses as Lemma~\ref{lemma:odd_iter_bound_theorem_4_1}, we have:
\begin{equation}
    \frac{E(j, n)}{n} \le \frac{\left(3 + \frac{1}{m}\right)^q - 3^q}{2^j}.
\end{equation}
\end{lemma}
\begin{proof}\leanok
By Lemma~\ref{lemma:odd_iter_bound_theorem_4_1}:
\begin{equation}
    T^j(n) \le \frac{n(3+1/m)^q}{2^j}.
\end{equation}
Substituting the factored form $T^j(n) = C(j,n) \cdot n + E(j,n)$ from Lemma~\ref{lemma:linear_decomposition'} and dividing through by $n > 0$:
\begin{equation}
    C(j,n) + \frac{E(j,n)}{n} \le \frac{(3+1/m)^q}{2^j}.
\end{equation}
Since $C(j, n) = 3^q / 2^j$, rearranging gives the result.
\end{proof}


  Since the harmonic mean $h$ is not less than the minimum $m$ of the odd iterates of $T$,
  Lemma~\ref{lemma:correction_ratio_bound_theorem_4_1} is a weaker version of
  Theorem~$4.1$ in Rozier--Terracol.


\begin{lemma}\label{lemma:paradoxical_ratio_real}\lean{CollatzMapBasics.paradoxical_ratio_real}\leanok
Let $j, q, m \in \mathbb{N}$ with $j > 0$ and $m > 0$. Suppose the following inequalities hold for the rationals $3^q/2^j$ and $(3+1/m)^q/2^j$:
\begin{equation}
    1 - \frac{3^q}{2^j} > 0 \quad \text{and} \quad 1 - \frac{3^q}{2^j} \le \frac{(3 + 1/m)^q - 3^q}{2^j}.
\end{equation}
Then the ratio $q/j$ satisfies:
\begin{equation}
    \frac{\log 2}{\log(3 + 1/m)} \le \frac{q}{j} < \frac{\log 2}{\log 3}.
\end{equation}
\end{lemma}
\begin{proof}\leanok
For the upper bound, $1 - 3^q/2^j > 0$ implies $3^q < 2^j$. Taking logs gives $q \log 3 < j \log 2$, so $q/j < \log 2 / \log 3$.
For the lower bound, the second inequality implies $2^j - 3^q \le (3+1/m)^q - 3^q$, which simplifies to $2^j \le (3+1/m)^q$. Taking logs gives $j \log 2 \le q \log(3+1/m)$, yielding the final result.
\end{proof}

\begin{corollary}\label{lemma:paradoxical_ratio_bounds}\lean{CollatzMapBasics.paradoxical_ratio_bounds}\uses{lemma:isParadoxical_bound,lemma:correction_ratio_bound_theorem_4_1,lemma:paradoxical_ratio_real}\leanok
Let $(j, n)$ be a paradoxical sequence starting at $n > 0$. Let $m > 0$ be a lower bound for all odd iterates $T^i(n)$ with $i < j$. Let $q$ be the number of odd steps. Then:
\begin{equation}
    \frac{\log 2}{\log(3 + 1/m)} \le \frac{q}{j} < \frac{\log 2}{\log 3}.
\end{equation}
\end{corollary}
\begin{proof}\leanok
This follows from Lemma~\ref{lemma:correction_ratio_bound_theorem_4_1} and the algebraic properties derived in Lemma~\ref{lemma:paradoxical_ratio_real}.
\end{proof}

  Since the harmonic mean $h$ is not less than the minimum $m$ of the odd iterates of $T$,
  Lemma~\ref{lemma:paradoxical_ratio_bounds} is a weaker version of Corollary~$4.2$ in Rozier--Terracol.

\begin{lemma}\label{lemma:log_ratio_not_int}\lean{CollatzMapBasics.log_ratio_not_int}\uses{lemma:irrational_xi}\leanok
For any $j \ge 2$ and any $k \in \mathbb{N}$, the value $j \cdot \log 2 / \log 3$ is not an integer:
\begin{equation}
    \frac{j \log 2}{\log 3} \neq k.
\end{equation}
\end{lemma}
\begin{proof}\leanok
If $j \log 2 / \log 3 = k$, then $\xi = \log 2 / \log 3 = k/j \in \mathbb{Q}$, contradicting the irrationality of $\xi$ (Lemma~\ref{lemma:irrational_xi}).
\end{proof}

\begin{lemma}\label{lemma:strict_floor_log_ratio}\lean{CollatzMapBasics.strict_floor_log_ratio}\uses{lemma:log_ratio_not_int}\leanok
For any $j \ge 2$, the floor of $j \log 2 / \log 3$ is strictly less than $j \log 2 / \log 3$:
\begin{equation}
    \left\lfloor \frac{j \log 2}{\log 3} \right\rfloor < \frac{j \log 2}{\log 3}.
\end{equation}
\end{lemma}
\begin{proof}\leanok
By definition, $\lfloor x \rfloor \le x$. The floor must be strictly less because $j \log 2 / \log 3$ is not an integer by Lemma~\ref{lemma:log_ratio_not_int}.
\end{proof}

\begin{lemma}\label{lemma:paradoxical_m_bound_q_real}\lean{CollatzMapBasics.paradoxical_m_bound_q_real}\leanok
Let $j, q, m \in \mathbb{N}$ with $j, q, m > 0$. Suppose the ratio bounds
\begin{equation}
    \frac{\log 2}{\log(3 + 1/m)} \le \frac{q}{j} < \frac{\log 2}{\log 3}
\end{equation}
hold. Then:
\begin{enumerate}
    \item $2^{j/q} - 3 > 0$, i.e.,\  $2^{j/q} > 3$;
    \item $m \le \dfrac{1}{2^{j/q} - 3}$.
\end{enumerate}
\end{lemma}
\begin{proof}\leanok
The upper bound $q/j < \log 2 / \log 3$ implies $q \log 3 < j \log 2$, hence $\log 3 < (j/q) \log 2 = \log(2^{j/q})$, so $2^{j/q} > 3$.

The lower bound $\log 2 / \log(3+1/m) \le q/j$ implies $j \log 2 \le q \log(3+1/m)$, i.e., $(j/q)\log 2 \le \log(3+1/m)$, so $2^{j/q} \le 3 + 1/m$. This gives $2^{j/q} - 3 \le 1/m$, and since $2^{j/q} - 3 > 0$, we can take reciprocals (reversing the inequality) to get $m \le 1/(2^{j/q} - 3)$.
\end{proof}

\begin{lemma}\label{lemma:paradoxical_m_bound_floor_real}\lean{CollatzMapBasics.paradoxical_m_bound_floor_real}\uses{lemma:paradoxical_m_bound_q_real,lemma:strict_floor_log_ratio}\leanok
Let $j, q, m \in \mathbb{N}$ with $j \ge 2$ and $m > 0$. Suppose the ratio bounds of Lemma~\ref{lemma:paradoxical_ratio_real} hold for $q/j$. Define:
\begin{equation}
    r = \frac{j}{\left\lfloor j \log 2 / \log 3 \right\rfloor}.
\end{equation}
Then $2^r - 3 > 0$ and $m \le 1/(2^r - 3)$.
\end{lemma}
\begin{proof}\leanok
From $q \le \lfloor j \log 2/\log 3 \rfloor$ (by the upper bound $q < j \log 2/\log 3$), we have $r = j/\lfloor j \log 2/\log 3 \rfloor \le j/q$.
This means $2^r \le 2^{j/q}$, so $1/(2^r - 3) \ge 1/(2^{j/q} - 3)$.
By Lemma~\ref{lemma:paradoxical_m_bound_q_real}, $m \le 1/(2^{j/q}-3) \le 1/(2^r - 3)$.

To show $2^r > 3$: since $\lfloor j\log 2/\log 3 \rfloor < j \log 2/\log 3$ (strictly, by Lemma~\ref{lemma:strict_floor_log_ratio}), $r = j/\lfloor j\log 2/\log 3 \rfloor > j/(j\log 2/\log 3) = \log 3/\log 2$.  Therefore $r \log 2 > \log 3$, i.e., $\log(2^r) > \log 3$, so $2^r > 3$.
\end{proof}

\begin{corollary}\label{corollary:paradoxical_m_bound}\lean{CollatzMapBasics.paradoxical_m_bound}\uses{lemma:paradoxical_ratio_bounds,lemma:paradoxical_m_bound_floor_real}\leanok
Let $(j, n)$ be a paradoxical pair with $j \ge 2$ and $n > 0$. Let $m > 0$ be a lower bound for all odd iterates $T^k(n)$ with $k < j$. Define $r = j / \lfloor j \log 2 / \log 3 \rfloor$. Then:
\begin{equation}
    2^r - 3 > 0 \quad \text{and} \quad m \le \frac{1}{2^r - 3}.
\end{equation}
\end{corollary}
\begin{proof}\leanok
Combine the ratio bounds from Lemma~\ref{lemma:paradoxical_ratio_bounds} with Lemma~\ref{lemma:paradoxical_m_bound_floor_real}.
\end{proof}

  Since the harmonic mean $h$ is not less than the minimum $m$ of the odd iterates of $T$,
  Lemma~\ref{corollary:paradoxical_m_bound} is weaker than Corollary~$4.3$ in Rozier--Terracol.
